\subsection{Conditional Probability}
\hrulefill

A conditional probability of an event is the probability that the event occurs given that some other event
has already occurred. Conditional probabilities can prove to be counterintuitive and may lead tableofcontents
suprising results.

\paragraph*{Definition:} The conditional probability of an event $A$ given that an event $B$ has occurred is denoted by
$P(A|B)$ and is defined as:
\[
    P(A|B) = \frac{P(A \cap B)}{P(B)}
\]

\paragraph*{Example:} If you flip 3 coins, what is the probability that exactly 2 are heads (B)?

\begin{align*}
    \Omega &= \{HHH, HHT, HTH, HTT, THH, THT, TTH, TTT\} \\
    B &= \{HHT, HTH, THH\} \\
    A &= \{HHH, HTT, HTH, HHT\} \\
\end{align*}
\begin{align*}
    P(A \cap B) &= \frac{2}{8} \\
    P(B) &= \frac{3}{8} \\
    P(A|B) = \frac{P(A \cap B)}{P(B)} &= \frac{\frac{2}{8}}{\frac{3}{8}} = \frac{2}{3}
\end{align*}

\subsubsection*{Multiplication Rule}
The multiplication rule is a way to find the probability of the intersection of two events. It states that:
\begin{align*}
    P(A \cap B) &= P(A|B) \cdot P(B) \\
    P(A \cap B) &= P(B|A) \cdot P(A) \\
\end{align*}

\pagebreak

CHECK THIS FOR ACCURACY
\paragraph*{General Multiplication Rule:} For any number of events $A_1, A_2, \ldots, A_n$:
\begin{align*}
    P\left(\bigcap_{i=1}^{n} A_i\right) &= P(A_1) \cdot P(A_2|A_1) \cdot P(A_3|A_1 \cap A_2) \cdots P(A_n|A_1 \cap A_2 \cap \cdots \cap A_{n-1}) \\
\end{align*}

\subsubsection*{Law of Total Probability}
For any events A and B, we can get the following from the multiplication rule:
\begin{align*}
    P(A \cap B) &= P(A|B)P(B) \\
    P(A \cap B') &= P(A|B')P(B') \\
\end{align*}:
Note that $B$ and $B'$ form a \textit{partition} of $\Omega$.

\paragraph*{General Law of Total Probability:}
Let $B_1, \ldots, B_n$ be a partition of the sample space $\Omega$. Then for any event $A$:
\[
    P(A) = \sum_{i=1}^{n} P(A|B_i)P(B_i)
\]

\paragraph*{Example:}
Suppose that we have 2 boxes with marbles in them. Box 1 has 1 white and 3 black marbles. Box 2 
has 2 white and 2 black marbles. We randomly select a box and then randomly select a marble from that box.
What is the probability that we select a black marble?

\begin{align*}
    P(Black) &= P(Black|Box 1)P(Box 1) + P(Black|Box 2)P(Box 2) \\
    &= \left(\frac{3}{4}\right)\left(\frac{1}{2}\right) + \left(\frac{2}{4}\right)\left(\frac{1}{2}\right) \\
    &= \frac{3}{8} + \frac{2}{8} = \frac{5}{8} \\
\end{align*}


\paragraph*{Example:}
On any given morning, it is either raining (25\% of the time), foggy (35\% of the time), or sunny (40\% of the time).
If it is raining, there is a 30\% change that I will miss the bus. If it is foggy, there is a 20\%
chance that I will miss the bus. If it is sunny, there is a 10\% chance that I will miss the bus.
What is the probability that I will miss the bus on any given morning?

\begin{align*}
    P(MB) &= P(MB|R)P(R) + P(MB|F)P(F) + P(MB|S)P(S) \\
    &= (0.30)(0.25) + (0.20)(0.35) + (0.10)(0.40) \\
    &= 0.075 + 0.07 + 0.04 = 0.185 \\
\end{align*}

\subsubsection*{Bayesian Analysis}
Bayesian analysis is a statistical method that involves using Bayes' theorem to update the probability of a hypothesis
based on new evidence or data. It is named after Thomas Bayes, an 18th-century statistician and theologian.

\paragraph*{Bayes' Theorem}
Let $B_1,\ldots,B_n$ form a partition of the sample space $\Omega$. Then for any event $A$:
\[
    P(B_i|A) = \frac{P(A|B_i)P(B_i)}{\sum_{j=1}^{n} P(A|B_j)P(B_j)}
\]
Note that in the special case of the partition consisting of $B$ and $B'$, this reduces to:

\[
    P(B|A) = \frac{P(A|B)P(B)}{P(A|B)P(B) + P(A|B')P(B')}
\]

\paragraph*{Example:}
If I pull a black marble from one of the boxes in the previous example, what is the probability that I pulled it from Box 1?
\begin{align*}
    P(B_1|Black) &= \frac{P(Black|B_1)P(B_1)}{P(Black)} \\
    &= \frac{(\frac{3}{4})(\frac{1}{2})}{\frac{5}{8}} = \frac{3}{5}
\end{align*}

\paragraph*{Example:}
According to the Arizona Chapter of the American Lung Association, 7\% of the population has lung disease. Of those,
people, 90\% are smokers; and of those without lung disease, 74.7\% are non-smokers. What are the chances that a smoker has 
lung disease?
\begin{align*}
    P(LD|S) &= \frac{P(S|LD)P(LD)}{P(S|LD)P(LD) + P(S|LD')P(LD')} \\
    &= \frac{(0.90)(0.07)}{(0.90)(0.07) + (0.253)(0.93)} \\
    &= \frac{0.063}{0.063 + 0.23529} = \frac{0.063}{0.29829} \approx 0.2111 \\
\end{align*}

\paragraph*{Example:}
You give a friend a letter to mail. He forgets to mail it with probability 0.2. Given that he mails it, 
the Post Office delivers it with probability 0.9. Given that the letter was not delivered, what is the probability
that it was not mailed.

\begin{align*}
    P(M'|D') &= \frac{P(D'|M')P(M')}{P(D'|M')P(M') + P(D'|M)P(M)} \\
    &= \frac{(1)(0.2)}{(1)(0.2) + (0.1)(0.8)} \\
    &= \frac{0.2}{0.2 + 0.08} = \frac{0.2}{0.28} = \frac{5}{7} \\
\end{align*}